\documentclass[10pt,a4paper]{amsart}
\usepackage[margin=19mm]{geometry}
\usepackage{amsmath,amssymb,mathtools}
\usepackage[scr=rsfs,cal=euler]{mathalfa}
\usepackage{xfrac}
\usepackage[unicode,hidelinks,bookmarks=false]{hyperref}
\newcommand{\ie}{i.e.\ }
\newcommand{\cf}{cf.\ }
\newcommand{\resp}{respectively\ }
\newcommand{\Subsection}{\subsection}
\providecommand{\nobreakdash}{\nobreak\hbox{-}}
\setlength{\emergencystretch}{2em}
\multlinegap0pt
\usepackage{tikz}
\usepackage{float}
\usepackage{amssymb}
\usepackage{mathtools}
\usepackage{bm}
\usepackage[shortlabels]{enumitem}
\usepackage{xcolor}
\newtheorem{proposition}{Proposition}
\newtheorem{theorem}{Theorem}
\usepackage{cleveref}
\crefname{proposition}{Proposition}{Propositions}
\crefname{theorem}{Theorem}{Theorems}
\newcommand{\Id}{\mathrm{Id}}
\newcommand{\Q}{Q}
\newcommand{\Qe}{Q_\varepsilon}
\newcommand{\eps}{\varepsilon}
\newcommand{\mc}[1]{\mathcal{#1}}
\title{\bf Thin Domains, Reduction, and Slow Manifolds}
\author{Christian Kuehn, Jan-Eric Sulzbach}
\date{Source: arXiv:2606.05398v1 (CC BY 4.0)}
\begin{document}
\maketitle

\noindent\emph{Source and licence.} \bf Thin Domains, Reduction, and Slow Manifolds. Version arXiv:2606.05398v1 (CC BY 4.0), distributed by arXiv under the Creative Commons Attribution 4.0 International licence, http://creativecommons.org/licenses/by/4.0/, as recorded in the arXiv licence metadata for this exact version and shown on the abstract page https://arxiv.org/abs/2606.05398v1. Full licence notice: \texttt{licenses/arxiv-2606.05398-CC-BY-4.0.txt}.\par\smallskip

\noindent\textit{Editorial scope.} Counted unit: the article's theorem prop:general-splitting (source0.tex lines 902--953). The article's complete original proof of it is delivered with it (source0.tex lines 956--1215, one \texttt{proof} environment of 260 lines). No mathematics is written, repaired or re-derived: the statement and the proof are the article's own lines, in the article's own notation, and no retained line is substituted or re-flowed.  Retained uncounted as the source-local context the proof needs: lines 754--757; lines 768--795; lines 781--784.  Nothing was omitted from any retained span, so the package carries no omission record.  No retained line contains a citation, so the wrapper carries no bibliography and no bibliography entry.  No retained line contains a figure environment, an image-inclusion command or drawing code, so no image is delivered and the package declares no asset dependency.  Editorial formatting only, in addition: an AMS \texttt{amsart} preamble that restates the article's own declarations and macros used by the retained lines, namely the environments \texttt{proposition}, \texttt{theorem} on separate counters, together with the standard packages those lines need.  The retained environments print in this document as Proposition 1, Theorem 1 in order of appearance, whereas the article prints the same items with the numbering of its own full document.  The complete native source of the article as distributed in the arXiv e-print for this exact version is bundled unchanged as \texttt{sources/202113/source0.tex}.
\begin{equation}\label{eq:general-physical}
\partial_t \hat u = \Delta \hat u + f(\hat u)
\qquad\text{in }\Qe,
\end{equation}

\begin{proposition}[Rescaled equation and conormal boundary conditions]\label{prop: rescaled thin domain}
 Let $\widehat u$ be a sufficiently smooth solution of \eqref{eq:general-physical}, and define
\[
u(x,Y,t):=\widehat u(x,\varepsilon g(x)Y,t).
\]
Set
\[
\mathcal D:=\partial_x-b(x,Y)\partial_Y,
\qquad
b(x,Y):=Y\frac{g'(x)}{g(x)} .
\]
Here and below \(\mc D^2\) denotes the composition \(\mc D\circ \mc D\).
Then, the transformed equation is given by
\begin{equation}\label{eq:rescaled-operator-general}
\partial_t u = \frac{1}{\eps^2 g(x)^2}\partial_Y^2u + \mathcal D^2 u + f(u)
\qquad\text{in }\Q.
\end{equation}
The pulled-back homogeneous Neumann conditions are
\begin{align*}
\partial_Yu=0
\qquad\text{on } \Omega \times \{0\},\\
\partial_Yu-\varepsilon^2g(x)g'(x)\mathcal D u=0
\qquad\text{on } \Omega\times \{1\},
\intertext{and}
\mathcal D u=0
\qquad\text{on }\partial\Omega\times(0,1).
\end{align*}  
\end{proposition}

\begin{equation}
\partial_t u = \frac{1}{\eps^2 g(x)^2}\partial_Y^2u + \mathcal D^2 u + f(u)
\qquad\text{in }\Q.
\end{equation}

\begin{theorem}[Fast-slow splitting on a variable-thickness thin strip]
\label{prop:general-splitting}
Let \(u\) be a sufficiently smooth solution of \eqref{eq:rescaled-operator-general} satisfying the pulled-back boundary conditions from \Cref{prop: rescaled thin domain}. Define
\[
v:=Mu,
\qquad
w:=(\Id-M)u,
\qquad
u=v+w.
\]
Then \(v\) is independent of \(Y\), \(Mw=0\), and all traces below are understood in the classical sense.
Set
\begin{align}
     B_g v&:=\frac1g\partial_x(g\partial_xv),\\
     \intertext{and} \label{eq: boundary term Rg}
     R_gw&:=-\frac1g\partial_x\bigl(g' w(\cdot,1)\bigr).
\end{align}
Then the rescaled problem can be written as
\begin{align}
\partial_t v
&=
 B_gv+R_gw+Mf(v+w),
\label{eq:slow-general}
\\
\partial_t w
&=
\frac1{\varepsilon^2g(x)^2}\partial_Y^2w
+\mathcal D^2w
-\mathcal E\bigl(\frac{g'(x)}{g(x)}\,\partial_xv+R_gw\bigr)
+(\Id-M)f(v+w),
\label{eq:fast-general}
\end{align}
where \(\mathcal E\phi(x,Y):=\phi(x)\).
The boundary conditions for the split variables are
\begin{align*}
\partial_Yw=0
\qquad\text{on } \Omega \times \{0\},\\
\partial_Yw-\varepsilon^2gg'\mathcal Dw
=
\varepsilon^2gg'\partial_xv
\qquad\text{on } \Omega\times \{1\},
\intertext{and}
\mathcal Dw=-\partial_xv
\qquad\text{on }\partial\Omega\times(0,1).
\end{align*}
In addition, averaging the lateral boundary condition gives the compatibility condition
\[
\partial_x v-\frac{g'}g w(\cdot,1)=0
\qquad\text{on }\partial\Omega .
\]
Moreover, $B_g=\partial_x^2$ and $R_g=0$ whenever \(g\) is constant.
\end{theorem}

\begin{proof}
We start with the averaged equation. 
The main point is that, unlike in the flat coordinate-Neumann case, the average of the singular transverse term produces a boundary contribution. Indeed,
\[
M\left(\frac1{\varepsilon^2g^2}\partial_Y^2u\right)
=
\frac1{\varepsilon^2g^2}
\int_0^1\partial_Y^2u\,dY
=
\frac1{\varepsilon^2g^2}
\bigl[\partial_Yu\bigr]_{Y=0}^{Y=1}.
\]
Using the transformed Neumann conditions,
\[
\partial_Yu|_{Y=0}=0,
\qquad
\partial_Yu|_{Y=1}
=
\varepsilon^2gg'\mathcal Du|_{Y=1},
\]
we obtain
\[
M\left(\frac1{\varepsilon^2g^2}\partial_Y^2u\right)
=
\frac{g'}g\,\mathcal Du|_{Y=1}.
\]
Next we compute \(M\mathcal D^2u\). For any sufficiently smooth \(\phi\), one has
\[
M(\mathcal D\phi)
=
\int_0^1(\partial_x\phi-Y \frac{g'}g\partial_Y\phi)\,dY.
\]
Since \(\frac{g'}g\) is independent of \(Y\),
\[
M(\mathcal D\phi)
=
\partial_xM\phi
-
\frac{g'}g\int_0^1Y\partial_Y\phi\,dY.
\]
Integration by parts in \(Y\) gives
\[
\int_0^1Y\partial_Y\phi\,dY
=
\bigl[Y\phi\bigr]_{0}^{1}
-
\int_0^1\phi\,dY
=
\phi(\cdot,1)-M\phi.
\]
Therefore
\[
M(\mathcal D\phi)
=
\partial_x M\phi
-
\frac{g'}g\phi(\cdot,1)
+
\frac{g'}g M\phi.
\]
Applying this first with \(\phi=u=v+w\), and using \(Mu=v\), \(Mw=0\), gives
\[
M(\mathcal Du)
=
\partial_xv-\frac{g'}g(v+w(\cdot,1))+\frac{g'}g v
=
\partial_xv-\frac{g'}g w(\cdot,1).
\]
Set
\[
P:=M(\mathcal Du)=\partial_xv-\frac{g'}g w(\cdot,1).
\]
Applying the same identity with \(\phi=\mathcal Du\), we get
\[
M(\mathcal D^2u)
=
\partial_xP
-
\frac{g'}g(\mathcal Du)(\cdot,1)
+
\frac{g'}gP.
\]
Substituting \(P=\partial_xv-\frac{g'}gw(\cdot,1)\) yields
\[
\begin{aligned}
M(\mathcal D^2u)
&=
\partial_x\bigl(\partial_xv-\frac{g'}gw(\cdot,1)\bigr)
-
\frac{g'}g(\mathcal Du)(\cdot,1)
+
\frac{g'}g\bigl(\partial_xv-\frac{g'}g w(\cdot,1)\bigr)
\\
&=
\partial_x^2v
+
\frac{g'}g\partial_xv
-
\partial_x\Bigl(\frac{g'}gw(\cdot,1)\Bigr)
-
\Bigl(\frac{g'}g\Bigr)^2w(\cdot,1)
-
\frac{g'}g(\mathcal Du)(\cdot,1).
\end{aligned}
\]
Adding the averaged singular term yields
\[
\begin{aligned}
M\left(
\frac1{\varepsilon^2g^2}\partial_Y^2u+\mathcal D^2u
\right)
&=
\frac{g'}g(\mathcal Du)(\cdot,1)
+
M(\mathcal D^2u)
\\
&=
\partial_x^2v
+
\frac{g'}g\partial_xv
-
\partial_x\Bigl(\frac{g'}g w(\cdot,1)\Bigr)
-
\Bigl( \frac{g'}g\Bigr)^2w(\cdot,1).
\end{aligned}
\]
The boundary contribution \(\frac{g'}g(\mathcal Du)(\cdot,1)\) cancels exactly. Thus, defining
\[
 B_gv:=\frac1g\partial_x(g\partial_xv),
\]
and
\[
R_gw
:=
-\partial_x\Bigl( \frac{g'}g w(\cdot,1)\Bigr)-\Bigl(\frac{g'}g\Bigr)^2w(\cdot,1)= -\frac1g\partial_x\bigl(g'w(\cdot,1)\bigr),
\]
we have
\[
M\left(
\frac1{\varepsilon^2g^2}\partial_Y^2u+\mathcal D^2u
\right)
=
B_gv+R_gw.
\]
Applying \(M\) to the equation therefore gives
\[
\partial_tv
=
 B_gv+R_gw+Mf(v+w).
\]

It remains to derive the equation for \(w\). Since \(w=u-v\), we have
\[
\partial_tw=\partial_tu-\partial_tv.
\]
Using
\[
\mathcal Dv=\partial_xv,
\qquad
\mathcal D^2v=\partial_x^2v,
\quad \text{and}\quad 
\partial_Yv=0,
\]
we obtain
\[
\begin{aligned}
\partial_tw
&=
\frac1{\varepsilon^2g^2}\partial_Y^2w
+
\mathcal D^2w
+
\partial_x^2v
+
f(v+w)
-
\left(
 B_gv+R_gw+Mf(v+w)
\right)
\\
&=
\frac1{\varepsilon^2g^2}\partial_Y^2w
+
\mathcal D^2w
-
\frac{g'}g\partial_xv
-
R_gw
+
(\Id-M)f(v+w).
\end{aligned}
\]
Viewing the \(x\)-dependent terms as functions on \(Q\) through the embedding
\[
(\mathcal E\phi)(x,Y):=\phi(x),
\]
this becomes
\[
\partial_tw
=
\frac1{\varepsilon^2g^2}\partial_Y^2w
+
\mathcal D^2w
-
\mathcal E\bigl(\frac{g'}g\partial_xv+R_gw\bigr)
+
(\Id-M)f(v+w).
\]

Finally, the boundary conditions for \(v\) and \(w\) follow by substituting \(u=v+w\) into the
transformed Neumann conditions. Since \(v\) is independent of \(Y\),
\[
\partial_Yu=\partial_Yw,
\qquad
\mathcal Du=\partial_xv+\mathcal Dw.
\]
Thus the lower boundary condition gives
\[
\partial_Yw=0
\qquad\text{on }Y=0.
\]
The upper boundary condition gives
\[
\partial_Yw-\varepsilon^2gg'(\partial_xv+\mathcal Dw)=0,
\]
or equivalently
\[
\partial_Yw-\varepsilon^2gg'\mathcal Dw
=
\varepsilon^2gg'\partial_xv
\qquad\text{on }Y=1.
\]
The lateral condition gives
\[
\partial_xv+\mathcal Dw=0
\qquad\text{on }\partial\Omega\times(0,1),
\]
that is,
\[
\mathcal Dw=-\partial_xv
\qquad\text{on }\partial\Omega\times(0,1).
\]
Averaging this last condition in \(Y\), and using
\[
M(\mathcal Dw)
=
-\frac{g'}gw(\cdot,1),
\]
gives the induced slow boundary condition
\[
\partial_xv- \frac{g'}gw(\cdot,1)=0
\qquad\text{on }\partial\Omega.
\]

If \(g\) is constant, then \(g'=0\), \(b=0\), and \(R_g=0\). Hence
\[
 B_g=\partial_x^2,
\]
and the system reduces to the standard flat-domain fast-slow splitting. This completes the proof.
\end{proof}

\end{document}
